\documentclass[10pt,a4paper]{amsart}
\usepackage[margin=19mm]{geometry}
\usepackage{amsmath,amssymb,mathtools}
\usepackage[scr=rsfs,cal=euler]{mathalfa}
\usepackage{xfrac}
\usepackage[unicode,hidelinks,bookmarks=false]{hyperref}
\newcommand{\ie}{i.e.\ }
\newcommand{\cf}{cf.\ }
\newcommand{\resp}{respectively\ }
\newcommand{\Subsection}{\subsection}
\providecommand{\nobreakdash}{\nobreak\hbox{-}}
\setlength{\emergencystretch}{2em}
\multlinegap0pt
\usepackage{bm}
\usepackage{bbm}
\usepackage{dsfont}
\usepackage{xspace}
\usepackage{cancel}
\usepackage{mathtools}
\usepackage{amsthm}
\makeatletter
\@ifundefined{theorem}{\newtheorem{theorem}{Theorem}}{}
\@ifundefined{lemma}{\newtheorem{lemma}[theorem]{Lemma}}{}
\makeatother
\newcommand{\Fix}{\operatorname{Fix}}
\title[A Hard-Core Subshift Whose Sofic Mean Dimension Depends on t]{A Hard-Core Subshift Whose Sofic Mean Dimension Depends on the Sofic Approximation}
\author{Xianqiang Li, Zhuowei Liu}
\date{Source: arXiv:2607.21398v1 (CC BY 4.0)}
\begin{document}
\maketitle

\noindent\emph{Source and licence.} A Hard-Core Subshift Whose Sofic Mean Dimension Depends on the Sofic Approximation. Version arXiv:2607.21398v1 (CC BY 4.0), distributed under the Creative Commons Attribution 4.0 International licence, as recorded in the arXiv licence metadata for this exact version and shown on the abstract page https://arxiv.org/abs/2607.21398v1. Full rights notice: licenses/arxiv-2607.21398-CC-BY-4.0.txt.\par\smallskip

\noindent\textit{Editorial scope.} Counted unit: the Lemma whose source label is \texttt{lem:random-free} in the article (source0.tex lines 729--739, source label \texttt{lem:random-free}), delivered together with the article's own complete original proof of it (source0.tex lines 741--835, one \texttt{proof} environment of 95 lines). No mathematics is written, repaired or re-derived: statement and proof are the article's own text line for line. Required context retained uncounted: none. Every definition, display and label of those lines is retained unchanged. Omitted from this package and not redistributed: 1--728, 740--740, 836--946. Nothing inside a retained excerpt is omitted or altered: every retained body is delivered exactly as the article's lines are written, so no omission receipt of any kind accompanies this package. Citations: the retained excerpts carry no citation command at all, so no bibliography entry of the article is transcribed and no reference is redistributed. Reference adaptation: none was needed and none was made; every reference inside the delivered unit resolves inside it, so no unresolved reference or citation is produced. Printed slips kept literal: no printed word, symbol, hypothesis or number of the counted statement or of its proof was altered. Layout and class: the article's own class is \texttt{amsart} with packages including lineno, geometry, newtxtext, newtxmath, mathrsfs, fancyhdr, url, xy, hyperref, enumitem; the wrapper is the lane's pinned \texttt{amsart} editorial wrapper at 10pt on A4, which restates the theorem-like environments the retained text needs and adds the article's own macro definitions that the retained text needs (\texttt{\textbackslash Fix}), with extra wrapper packages (bm, bbm, dsfont, xspace, cancel, mathtools). The delivered numbering is the unit's own. Assets: no retained span contains a figure environment, a graphics-inclusion command, a PDF-page inclusion, a table or a TikZ picture; no image file is copied into this package, the package declares no asset dependency and no image file is redistributed with it. Licence: version arXiv:2607.21398v1 of this article is distributed under the Creative Commons Attribution 4.0 International licence, as recorded in the arXiv licence metadata for this exact version and shown on the abstract page https://arxiv.org/abs/2607.21398v1. A full rights notice is delivered at licenses/arxiv-2607.21398-CC-BY-4.0.txt. Characters: every retained excerpt is pure ASCII, so the delivered engine, XeLaTeX with its default Latin Modern text font, reports no missing character.
	\begin{lemma}[Random permutations are asymptotically free]\label{lem:random-free}
		For every nontrivial reduced word $w\in F_2$ there is a constant $C_w<\infty$ such that
		\[
		\mathbb E|\Fix(\sigma(w))|\le C_w
		\]
		for all sufficiently large $d$.  Consequently, when $d\longrightarrow\infty$,
		\[
		\frac{|\Fix(\sigma(w))|}{d}\longrightarrow0
		\]
		in probability.
	\end{lemma}

	\begin{proof}
Write the reduced word as
\[
w=t_\ell\cdots t_1,
\qquad
 t_i\in\{a,a^{-1},b,b^{-1}\}.
\]
Fix a vertex $v\in[d]$ and follow the path described by the word $w$:
\[
x_0=v,
\qquad
x_i=\sigma(t_i)x_{i-1}
\quad(1\leq i\leq\ell).
\]
Then
\[
x_\ell=\sigma(w)v.
\]
Thus $v$ is fixed by $\sigma(w)$ exactly when $x_\ell=x_0$.

We reveal the values of $\pi_a$ and $\pi_b$ only when they are needed along
this path. Suppose that, before step $i$, the vertices
\[
x_0,x_1,\ldots,x_{i-1}
\]
are all distinct. Since $w$ is reduced, step $i$ cannot simply reverse the
edge used at step $i-1$. Hence the required value of $\pi_a$, $\pi_a^{-1}$,
$\pi_b$, or $\pi_b^{-1}$ has not yet been determined.

At most $i-1$ values of the relevant permutation have already been exposed.
Therefore the next vertex $x_i$ is uniformly distributed among at least
\[
d-(i-1)\geq d-\ell
\]
possible vertices. Among these possibilities, at most the $i$ previously
visited vertices
\[
x_0,\ldots,x_{i-1}
\]
produce a collision. Hence
\[
\mathbb P\bigl(\text{the first collision occurs at step }i\bigr)
\leq \frac{i}{d-\ell}.
\]

If $x_\ell=x_0$, then a collision must occur at some step
$1\leq i\leq\ell$. By the union bound,
\[
\begin{aligned}
\mathbb P\bigl(\sigma(w)v=v\bigr)
&\leq
\sum_{i=1}^{\ell}\frac{i}{d-\ell}\\
&=
\frac{\ell(\ell+1)}{2(d-\ell)}.
\end{aligned}
\]
For $d>2\ell$, this gives
\[
\mathbb P\bigl(\sigma(w)v=v\bigr)
\leq \frac{\ell(\ell+1)}{d}.
\]

Now sum over all $v\in[d]$. Since
\[
|\Fix(\sigma(w))|
=
\sum_{v=1}^{d}\mathbf 1_{\{\sigma(w)v=v\}},
\]
linearity of expectation gives
\[
\begin{aligned}
\mathbb E|\Fix(\sigma(w))|
&=
\sum_{v=1}^{d}\mathbb P\bigl(\sigma(w)v=v\bigr)\\
&\leq
 d\cdot\frac{\ell(\ell+1)}{d}\\
&=
\ell(\ell+1).
\end{aligned}
\]
Thus we may take $C_w=\ell(\ell+1)$.

Finally, for every $\eta>0$, by Markov's inequality, when $d \longrightarrow \infty$, 
\[
\mathbb P\left(
\frac{|\Fix(\sigma(w))|}{d}>\eta
\right)
\leq
\frac{\mathbb E|\Fix(\sigma(w))|}{\eta d}
\leq
\frac{\ell(\ell+1)}{\eta d}
\longrightarrow0.
\]
This proves the claimed convergence in probability.
\end{proof}

\end{document}
